\section{Introduction}
\label{sec-intro}

An ontology is a class hierarchy whose classes have names, and names can be
embedded: the Gene Ontology, the Human Phenotype Ontology, and ChEBI each give tens
or hundreds of thousands of terms a vector, so that terms close in meaning are close
in cosine similarity.
Tasks such as ontology matching, term normalization, and retrieval over annotated
data then ask for \emph{the terms nearest to a given term that lie under a given
class}: the nearest phenotype to ``short stature'' under ``abnormality of the
skeletal system'', the nearest process to ``apoptotic process'' under ``cell
death''.
The hierarchy defines the filter, and the vectors rank what passes it.

Vector databases answer such queries by giving the index a predicate, and they leave
four problems open.
\textbf{C1. A sorted index does not hold meaning.}
Triple stores answer patterns with range scans over sorted permutations, and it is
tempting to believe that terms close in meaning are close in the order, so that a
scan could serve similarity.
No linear order can keep all neighbors close (Section~\ref{sec-order}), and on four
ontologies the best order we tried finds 38\% to 56\% of a term's ten nearest
neighbors within 100 positions.
\textbf{C2. Filtered approximate search loses recall without warning.}
An IVF or graph index with a filter probes the lists or the walk nearest the query;
when few of the terms there pass the filter, it returns fewer, worse answers and
reports no error.
On our workload FAISS indexes with an ID selector return recall 0.61 (IVF) and 0.64
(HNSW), and 0.02 for HNSW when the query lies outside a small class.
\textbf{C3. The filter is correlated with the vectors.}
Filtered search is costed by assuming the filter is independent of the similarity,
so that $k/f$ candidates suffice for a filter that passes a fraction $f$.
In an ontology the neighbors of a term lie under the term's own ancestors, so the
estimate is wrong in both directions (Section~\ref{sec-filter}).
\textbf{C4. The filter lives in another system.}
A vector database knows nothing of the hierarchy: the caller must resolve a class to
its members, keep a table of them, and keep it current.

\begin{example}
\label{exa-motivation}
Let $q$ be \emph{apoptotic process} and $c$ the class \emph{cell death} in the Gene
Ontology.
A vector index returns neighbors from anywhere in the ontology, and a post-filter
keeps those under $c$ only if enough of them are.
For a query from outside a class of at most ten members, an ID-selector IVF returns
on average 5\% of the right answers.
Here the engine finds the 97 members of $c$ with one range scan of the B+ tree that
holds the triples, scores them exactly, and returns \emph{apoptotic process involved
in development} (cosine 0.865) first, in under a millisecond in a web page.
\end{example}

\stitle{Contributions}.
We build the filter into the store.
\begin{enumerate}
\item \textbf{A bound and a measurement on order} (Section~\ref{sec-order}): every
  linear order of a graph's vertices separates some edge by $(N{-}1)/D$ positions,
  and on four ontologies no order of the terms holds more than 56\% of the ten
  nearest neighbors within 100 positions.
\item \textbf{The hierarchy as a membership test} (Section~\ref{sec-filter}):
  postorder intervals in the engine's B+ tree make a class's members a few range
  scans and membership a binary search, which lets a vector search test the filter
  per row without resolving it.
\item \textbf{A filtering rule that does not lose recall} (Section~\ref{sec-filter}):
  four strategies, a measurement that the class filter is correlated with the
  vectors, a cost model for the crossover between resolving the class and testing
  it, and a rule that picks between them and is exact for the classes where
  approximate indexes fail.
\item \textbf{Vector indexes whose lists live in the triple tree}
  (Section~\ref{sec-index}): IVF posting lists stored as key ranges of the same
  B+ tree, so that they page, cache, and latch with the triples.
\end{enumerate}
We evaluate against FAISS, hnswlib, and DuckDB on four ontologies, with queries that are
stored terms and queries that are not (Section~\ref{sec-exp}).
The code, data scripts, and benchmark programs are at the artifact URL above.
